\chapter{06: Implementation and Integration with WanderWise+ Backend}

\section{1. WanderWise+ Integration Overview}

The geographic clustering module integrates with the WanderWise+ backend as a critical component in the tour planning pipeline. This integration connects the user preference processing stage (Module I NLC) with the route optimization stage (Module IV Genetic Algorithm), transforming a flat list of POIs into structured daily clusters ready for routing.

The clustering service receives POI data from the database layer, applies geographic clustering algorithms, and returns structured cluster assignments to the calling components. The service handles the complete workflow: retrieving POIs by category, validating user constraints (trip duration), performing K-means clustering with edge case handling, and formatting output for downstream processing.

The integration architecture positions clustering as a stateless service, enabling horizontal scaling to handle multiple concurrent requests. Each clustering request is independent, requiring no shared state between requests.

\section{2. Database Schema for Clustering}

The WanderWise+ database schema supports clustering operations through POI storage with geographic coordinates and a precomputed distance matrix:

\begin{itemize}
\item \textbf{goa\_places}: Stores POI coordinates essential for Haversine distance calculations
\item \textbf{distance\_matrix}: Provides precomputed distances for common POI pairs
\item \textbf{user\_trips}: Stores trip duration as a user preference
\item \textbf{itinerary\_clusters}: Stores clustering results for generated itineraries
\end{itemize}

\begin{sql}
-- POI table with geographic coordinates
CREATE TABLE goa_places (
    id UUID PRIMARY KEY DEFAULT gen_random_uuid(),
    name VARCHAR(255) NOT NULL,
    latitude DECIMAL(9, 6) NOT NULL,
    longitude DECIMAL(9, 6) NOT NULL,
    category VARCHAR(100) NOT NULL,
    popularity_score DECIMAL(3, 2) DEFAULT 5.0,
    avg_visit_duration_minutes INTEGER DEFAULT 60
);

-- Precomputed distance matrix
CREATE TABLE distance_matrix (
    poi1_id UUID REFERENCES goa_places(id),
    poi2_id UUID REFERENCES goa_places(id),
    distance_km DECIMAL(10, 2) NOT NULL,
    travel_time_minutes INTEGER,
    PRIMARY KEY (poi1_id, poi2_id)
);
\end{sql}

\section{3. K-Means Service Class}

The K-Means service class provides the production implementation of geographic clustering:

\begin{itemize}
\item \textbf{MIN\_POI\_PER\_DAY}: Minimum 6 POIs per cluster
\item \textbf{MAX\_POI\_PER\_DAY}: Maximum 15 POIs per cluster
\item \textbf{OUTLIER\_THRESHOLD\_KM}: 30 km for outlier detection
\item \textbf{DEFAULT\_RANDOM\_SEED}: 42 for reproducibility
\end{itemize}

The service class includes methods for:
\begin{itemize}
\item \textbf{cluster\_pois}: Main clustering interface
\item \textbf{\_detect\_outliers}: Identify geographic outliers
\item \textbf{\_perform\_clustering}: Execute K-means algorithm
\item \textbf{\_handle\_edge\_cases}: Manage imbalanced clusters
\end{itemize}

\section{4. API Endpoint Implementation}

The clustering service is exposed through a FastAPI endpoint:

\begin{itemize}
\item \textbf{Endpoint}: POST /api/clustering/geographic
\item \textbf{Input}: ClusteringRequest with poi\_ids, num\_days, optional random\_seed
\item \textbf{Output}: ClusteringResponse with clusters, warnings, iterations, converged
\end{itemize}

\begin{python}
@router.post("/api/clustering/geographic", response_model=ClusteringResponse)
def cluster_pois_geographic(request: ClusteringRequest):
    """
    Perform geographic clustering of POIs for multi-day tour planning.
    """
    session = get_db_session()
    pois = session.query(POI).filter(POI.id.in_(request.poi_ids)).all()
    service = GeographicClusteringService(db_session=session)
    result = service.cluster_pois(poi_objects, request.num_days, request.random_seed)
    return ClusteringResponse(success=True, clusters=cluster_responses, ...)
\end{python}

\section{5. Multi-Day Tour Workflow}

The complete multi-day tour workflow integrates clustering with the genetic algorithm:

\begin{enumerate}
\item \textbf{Classify}: Use NLC to filter POIs by user interests
\item \textbf{Cluster}: Apply geographic clustering with K = trip duration
\item \textbf{Optimize}: Use GA to optimize routes within each cluster
\end{enumerate}

\begin{python}
def generate_multi_day_itinerary(user_interests: List[str], 
                                 num_days: int) -> Dict:
    classifier = NLCClassifier()
    relevant_pois = classifier.filter_pois_by_interests(user_interests)
    
    clustering_service = GeographicClusteringService()
    cluster_result = clustering_service.cluster_pois(relevant_pois, num_days)
    
    optimizer = GARouteOptimizer()
    
    itinerary = {'days': []}
    for day_num, pois in cluster_result.clusters.items():
        route = optimizer.optimize_route(pois)
        itinerary['days'].append({'day_number': day_num + 1, 'pois': route})
    
    return itinerary
\end{python}

\section{6. Performance Considerations}

Geographic clustering performance depends on several factors:

\textbf{Computational Complexity}:
\begin{equation}
O(iterations \times K \times N)
\end{equation}

Where iterations is typically 10-20, K is the number of days (3-7), and N is the number of POIs (30-60). This results in approximately 2,000-8,000 distance calculations per clustering operation.

\textbf{Caching Strategies}:
\begin{itemize}
\item Precomputed distance matrix for POI-to-POI distances
\item Cached clustering results for common interest combinations
\end{itemize}

\textbf{Memory Usage}:
Minimal—clustering requires storing the POI list, cluster assignments, and centroids. For 100 POIs, this is less than 100KB of data.

The service can handle hundreds of concurrent clustering requests on standard server hardware.

\section*{References}

\begin{itemize}
\item Pedregosa, F., et al. (2011). Scikit-learn: Machine learning in Python. Journal of Machine Learning Research.
\item FastAPI Documentation. https://fastapi.tiangolo.com/
\item PostgreSQL PostGIS Documentation. https://postgis.net/
\item Arthur, D., \& Vassilvitskii, S. (2007). k-means++: The advantages of careful seeding.
\end{itemize}
